  \documentclass{article}
\usepackage{amsmath}
\usepackage{graphicx}
\graphicspath{ {images/} }
\usepackage[spanish]{babel}
\usepackage{bbm}
\usepackage{amsthm}
\usepackage{authblk}
\usepackage{hyperref}
\usepackage[utf8]{inputenc}
\usepackage{mathrsfs}
\usepackage{amsfonts}
\usepackage{amssymb}
\usepackage{enumerate}
\usepackage[left=3cm,right=2cm,top=2cm,bottom=2cm]{geometry}
\usepackage{epsfig}
 \renewcommand{\baselinestretch}{0.93}
 \thispagestyle{empty}
 \pagestyle{empty}
\setlength{\textheight}{53pc} \setlength{\textwidth} {36pc}
\setlength{\topmargin}{-4mm}
\usepackage{multicol}
\newcommand*{\email}[1]{%
    \normalsize\href{mailto:#1}{#1}\par
    }
\setlength{\columnsep}{1cm}
\newcommand{\N}{\mathbb{N}}
\newcommand{\calM}{\cal{M}}
\newcommand{\calP}{\cal{P}}
\newcommand{\F}{\mathbb{F}}
\newcommand{\R}{\mathbb{R}}
\newcommand{\Z}{\mathbb{Z}}
\newcommand{\Q}{\mathbb{Q}}
\newcommand{\C}{\mathbb{C}}
\newcommand{\bbH}{\mathbb{H}}


\theoremstyle{definition}
\newtheorem{ejer}{Ejercicio}
\author{Marcos Morales}
\affil{Universidad Finis Terrae}
\title{Semana 4\\}



	


\begin{document}


\begin{picture}(400, 75)
   \put(-40,15){\includegraphics[width=5cm]{UFT.png}}


\end{picture}


\begin{center}
\huge{Clase 6}
\end{center}

\begin{center}
\Large{ Marcos Morales, Camila Muñoz}
\end{center}

\begin{center}
	\large{Ecuaciones Diferenciales Ordinarias (EDO)}
\end{center}
\begin{center}
\setlength{\parindent}{0cm}\large{Clases centradas para capacitar a los estudiantes de ingeniería en Ecuaciones diferenciales ordinarias. \\}
\end{center}


\vspace{1cm}

\begin{center}
	\large{Contenidos:}
\end{center}

\setlength{\parindent}{2.95cm}- Ecuaciones  Lineales Homogéneas con coeficientes constantes \\
\hspace*{2.95cm}- Ejercicios para control 1.


\vspace{6cm}


\begin{center}
\today
\end{center}
\begin{center}
Santiago, Chile
\end{center}


\pagestyle{empty}


\setcounter{page}{1}
\pagestyle{plain}
\setlength{\parindent}{0cm}





\newpage





\newpage

\section{Ecuaciones  Lineales Homogéneas con coeficientes constantes}


Supongamos que tenemos la ecuación diferencial 

\begin{center}
$ay''+by'+cy=0$
\end{center}

donde $a,b,c \in \mathbb{R}$ $a \neq 0$ son constantes. Asumimos que la solución es de la forma $e^{mx}$ con $m \in\mathbb{R}$. Entonces reemplazando en la ecuación diferencial, tenemos que



\begin{center}
$am^2e^x+bme^{xm}+ce^{mx}=0$
\end{center}


Luego


\begin{center}
$e^{mx}(am^2+bm+c)=0$
\end{center}

Si tuviéramos que $am^2+bm+c=0$, entonces tendríamos que la función satisface la ecuación diferencial, entonces tenemos que resolver


\begin{center}
$am^2+bm+c =0$
\end{center}

Las soluciones son

\begin{center}
$\displaystyle m = \frac{-b \pm \sqrt{b^2-4ac}}{2a}$
\end{center}


Hay 3 casos a considerar \\


\textbf{Caso 1: } $b^2-4ac >0$ \\

En este caso obtenemos dos soluciones $m_1 = \displaystyle  \frac{-b + \sqrt{b^2-4ac}}{2a}$ y $\displaystyle m_2 = \frac{-b - \sqrt{b^2-4ac}}{2a}$. Las soluciones son 

\begin{center}
$y_1 = e^{m_1x}$ y $y_2 = e^{m_2x}$
\end{center}

Por lo tanto la solución general viene dada por

\begin{center}
$y=c_1e^{m_1x}+c_2e^{m_2 x}$
\end{center}


\textbf{Caso 2: } $b^2-4ac =0$ \\

En este caso obtenemos una solución $m_1 = \displaystyle  \frac{-b}{2a}$, por lo tanto una solución de la ecuación diferencial es $y_1= e^{m_1x}$. Para obtener la segunda solución usamos el método descrito en la sección pasada, en este caso la ecuación diferencial es

\begin{center}
$ay''+by'+cy=0$
\end{center}

Y dividiendo por $a$

\begin{center}
$\displaystyle y''+\frac{b}{a} y'+ \frac{c}{a} y=0$
\end{center}

Por lo tanto, en este caso $b(x)= \displaystyle \frac{b}{a}$. Reemplazando en la fórmula obtenemos


\begin{align*}
y_2 &= \displaystyle e^{m_1x} \int \frac{e^{-\int \frac{b}{a} dx}}{\left( e^{\frac{-b}{2a}x}\right)^2}dx \\
&= \displaystyle e^{m_1x} \int \frac{e^{-\frac{b}{a} x}}{e^{\frac{-b}{a}x}}dx \\
&= \displaystyle e^{m_1x} \int 1 dx \\
&= xe^{m_1x} \\
\end{align*}


Por lo tanto la solución general viene dada por


\begin{center}
$y=c_1e^{m_1x}+c_2xe^{m_1 x}$
\end{center}




\textbf{Caso 3: } $b^2-4ac < 0$ \\

En este caso obtenemos dos soluciones imaginarias distintas de la forma

\begin{center}
$m_1= \alpha + \beta i$
\end{center}

\begin{center}
$m_2= \alpha - \beta i$
\end{center}


En este caso las soluciones son imaginarias $e^{(\alpha + \beta i)x}$ y $e^{(\alpha - \beta i)x}$, pero queremos que sean reales, por lo tanto escribimos 


\begin{center}
$e^{(\alpha + \beta i)x} =e^{\alpha x}\cos(\beta x)+ e^{\alpha}\sen(\beta x)i$
\end{center}

\begin{center}
$e^{(\alpha - \beta i)x} = e^{\alpha x}\cos(\beta x) - e^{\alpha x}\sen(\beta x)i$
\end{center}

Podemos generar las nuevas soluciones


\begin{center}
$y_1 = \displaystyle \frac{e^{(\alpha + \beta i)x}+ e^{(\alpha - \beta i)x} }{2} = e^{\alpha x}\cos(\beta x)$
\end{center}

\begin{center}
$y_2 = \displaystyle \frac{e^{(\alpha + \beta i)x}- e^{(\alpha - \beta i)x} }{2i} = e^{\alpha x}\sen(\beta x)$
\end{center}

Que son soluciones reales. Por lo tanto la solución general viene dada por


\begin{center}
$y=c_1e^{\alpha x}\cos(\beta x)+c_2e^{\alpha x}\sen(\beta x)$
\end{center}



\textbf{Ejemplo 1:}  Resolver la ecuación diferencial $y'' + y' -6y =0$ \\

\textit{Respuesta: } sustituyendo por $e^{mx}$, obtenemos

\begin{center}
$m^2+m-6=0$
\end{center}

Las soluciones son $m_1=-3$ y $m_2 = 2$, estamos en el caso 1, por lo tanto la solución general es

\begin{center}
$y= c_1e^{-3x} + c_2e^{2x}$
\end{center}



\textbf{Ejemplo 2:}  Resolver la ecuación diferencial $y'' + w^2 y=0$ donde $w$ es una constante \\

\textit{Respuesta: } sustituyendo por $e^{mx}$, obtenemos

\begin{center}
$m^2+\omega^2=0$
\end{center}

Las soluciones son $m_1= \omega i$ y $m_2 = -\omega i$, estamos en el caso 3, donde $\alpha =0$ y $\beta = \omega$, las soluciones son

\begin{center}
$y_1 = e^{0x} \cos(\omega x) = \cos(\omega x)$
\end{center}

\begin{center}
$y_2 = e^{0x} \sen(\omega x) = \sen(\omega x)$
\end{center}

Por lo tanto, la solución general es

\begin{center}
$y= c_1\cos(\omega x) + c_2 \sen(\omega x)$
\end{center}



\newpage

\section{Ejercicios para control 1.}
\begin{enumerate}
\item Resuelva la siguiente ecuación diferencial de primer orden

\begin{align*}
  \frac{dy}{dx} &= -\frac{y\cos(x)+2x}{\sen(x)+\cos(y)}
\end{align*}


\textbf{Solución:} Reescribiendo la ecuación

\begin{align*}
  (y\cos(x)+2x) dx+(\sen(x)+\cos(y) )dy &= 0
\end{align*}
Derivando parcialmente
\begin{align*}
  \frac{\partial }{\partial x}(\sen(x)+\cos(y)) &= \cos(x) =\frac{\partial }{\partial y}(y\cos(x)+2x),
\end{align*}
verificamos que la ecuaci\'on es exacta. Integrando
\begin{align*}
  f(x,y)=\int y\cos(x)+2x dx&= y\sen(x)+x^2+c(y)
\end{align*}
para encontrar $c(y)$
\begin{align*}
  \frac{\partial }{\partial y} ( y\sen(x)+x^2+c(y))&=\sen(x)+c'(y)= \sen(x)+\cos(y)\\
  c'(y)&=\cos(y)\\
    c(y)&=\sen(y)\\
\end{align*}
Por lo que 
\begin{align*}
  f(x,y)= y\sen(x)+x^2+\sen(y)=c
\end{align*}
\item Resuelva la siguiente ecuación diferencial de primer orden
\[x^2y^2 dx+(x^3y+y+3)dy=0\]
\textbf{Solución:} Consideramos $M(x,y)=x^2y^2$ y $N(x,y)=x^3y+y+3$.

Como $M_y=2x^2y$ y $N_x=3x^2y$, tenemos que para obtener el factor integrante, debemos resolver la ecuación diferencial

    \[\frac{dw}{dy}=w\frac{N_x-M_y}{M}=\frac{w}{y},\]
por lo tanto,
\[\frac{dw}{w}=\frac{dy}{y}.\]
Integrando podemos considerar $w=y$.
Multilplicando la ecuación diferencial
\[x^2y^3 dx + (x^3y^2+y^2+3y)dy=0.\]

En este caso se verifica que $M_y=N_x=3x^2y^2$. Integrando 
\[\int M(x,y) dx=\frac{x^3y^3}{3}+c(y).\]
pero \[x^3y^2+c'(y)=\frac{\partial f}{\partial y}=N=x^3y^2+y^2+3y,\]
luego
\[c'(y)=y^2+3y.\]
Integrando y reemplazando
\[f(x,y)=\frac{x^3y^3}{3}+c(y)=\frac{x^3y^3}{3}+\frac{y^3}{3}+\frac{3}{2}y^2=c.\]
Si $y(0)=1$, 
\[f(0,1)=\frac{1^3}{3}+\frac{3}{2}1^2=\frac{11}{6}.\]



\item Resuelva la siguiente ecuación diferencial de primer orden
\[\frac{dy}{dx}=\frac{y}{2x}+\frac{x^2}{2y}\]
\textbf{Solución:} Reescribiendo, obntenemos una ecuación de Bernoulli con $n=-1$
\[\frac{dy}{dx}-\frac{1}{2x}y=\frac{x^2}{2}y^{-1},\]
Por lo tanto hacemos el cambio de variable $z=y^2$ y como $\dfrac{dz}{dx}=2yy'$ obtenemos la siguiente ecuación lineal
\[\frac{z'}{2}-\frac{z}{2x}=\frac{x^2}{2}\]
que es equivalente a 
\[z'-\frac{z}{x}=x^2\]
con $a(x)=-\dfrac{1}{x}$ y $b(x)=x^2$.
Usando directamente la fórmula
\begin{align*}
    z&=e^{-\int a(x)dx}\left(\int b(x)e^{\int a(x)dx}dx+c\right)\\
   z&=e^{\int 1/xdx}\left(\int x^2e^{\int -1/xdx}dx+c\right)\\
    z&=x\left(\int xdx+c\right)\\
    z&=x\left(\frac{x^2}{2}dx+cx\right)\\
    z&=\frac{x^3}{2}+cx^2\\
\end{align*}
Reemplanzando 
\[    y^2=\frac{x^3}{2}+cx^2.\]

Si $y(0)=1$, $ 1^2=\frac{0^3}{2}+c0^2$, entonces $0=1$ lo cual es imposible, por lo tanto no existe una soluci\'on con dicho valor inicial.

    \item Resuelva la siguiente ecuación diferencial de primer orden
        
     \[\frac{dy}{dx}=\frac{2x+2y}{2x+2y+2}=\frac{x+y}{x+y+1}\]   
\textbf{Solución:} Podemos hacer el cambio de variable $z=x+y$, luego $\dfrac{dz}{dx}=1+\dfrac{dy}{dx}$. Reemplazando
\begin{align*}
    \frac{dz}{dx}-1&=\frac{z}{z+1}\\
    \frac{dz}{dx}&=\frac{z}{z+1}+1\\
    \frac{dz}{dx}&=\frac{z}{z+1}+\frac{z+1}{z+1}=\frac{2z+1}{z+1}\\
\end{align*}
como es de variables separables tenemos que
\[\frac{z+1}{2z+1}dz=dx\]
Integrando a ambos lados haciendo el $u=2z+1$
\begin{align*}
    \frac{1}{4}\int \frac{u+1}{u}du=\int dx =x+c
\end{align*}
luego
\begin{align*}
    x+c&=\frac{1}{4}\left(u+\ln(u)\right)=\frac{1}{4}\left(2z+1+\ln(2z+1)\right)\\
        x+c&=\frac{1}{4}\left(2(x+y))+1+\ln(2(x+y)+1)\right).
\end{align*}
Si $y(0)=0$, tenemos que $0+c=\dfrac{1}{4}\left(2(0))+1+\ln(2(0)+1)\right)$, entonces $c=\dfrac{1}{4}\left(1+\ln(1)\right)=\dfrac{1}{4}$.


\end{enumerate}


\end{document} 
